\hypertarget{how-far-the-focal-plane-opa-scales}{%
\chapter{How far the focal-plane OPA
scales}\label{scaling-how-far-the-focal-plane-opa-scales}}

Follow-up to \texttt{docs/fpr\_design.md}. Code:
\texttt{scripts/m12\_scaling.py}; results \texttt{results/m12\_*.json}.
Packing is simulated on the beam-propagation matrix of design D1;
aberration with exact path lengths; the rest of the scaling table is
derived from those results. Values are SIM or DERIVED and valid in the
regime stated.

\hypertarget{1-result}{%
\section{Result}\label{scaling-1-result}}

One focal-plane engine moves about 40 independent tweezers along a row
at its present size (128 ports), and about 60 when the row is matched to
a +/-100 um objective field (192 ports). The packing limit is a spacing
of 3 um, i.e. 2.8 resolution cells or 1.75 spot widths. Below that the
spots merge. Power per tweezer does not grow with the number of
tweezers. A second dimension needs one engine per row: stacking chips
with fast-axis micro-lenses, as in laser-diode stacks, gives an
estimated 8 to 16 rows. Hence several hundred independent tweezers (16
rows x 40 to 60) are a realistic target, against a median of 13 (maximum
38) moved in parallel with crossed AODs in the 1,024-atom assembly cited
in the proposal.

The binding limits, in order: (i) the diffraction-limited field of the
objective, which sets the ports per row; (ii) the number, speed and
accuracy of fast actuators (2 per port, \textgreater= 200 kHz,
\textasciitilde1\% and 10 mrad); (iii) laser power through each chip;
(iv) slab size. Star-coupler aberration is not a limit below about 512
ports.

\hypertarget{2-tweezers-per-row-packing}{%
\section{Tweezers per row:
packing}\label{scaling-2-tweezers-per-row-packing}}

A +/-60 um row of design D1 filled at spacing s, neighbours in phase
quadrature, least-squares synthesis:

{\small
\begin{longtable}[]{@{}>{\raggedright\arraybackslash}p{0.092\dimexpr(\linewidth-10\tabcolsep)\relax}>{\raggedright\arraybackslash}p{0.101\dimexpr(\linewidth-10\tabcolsep)\relax}>{\raggedright\arraybackslash}p{0.294\dimexpr(\linewidth-10\tabcolsep)\relax}>{\raggedright\arraybackslash}p{0.303\dimexpr(\linewidth-10\tabcolsep)\relax}>{\raggedright\arraybackslash}p{0.210\dimexpr(\linewidth-10\tabcolsep)\relax}@{}}
\toprule\noalign{}
spacing & tweezers & intensity midway between neighbours &
others\textquotesingle{} change when one moves 0.25 um & residual after
switch-off \\
\midrule\noalign{}
\endhead
\bottomrule\noalign{}
\endlastfoot
6 um & 21 & 5e-4 & 0.11\%, 1.6 nm & \textless{} 1e-5 \\
4 um & 31 & 3.8\% & 0.10\%, 2.6 nm & 1e-5 \\
3 um & 41 & 22\% & 0.11\%, 3.3 nm & 9e-4 \\
2.5 um & 49 & 43\% & 0.79\%, 10 nm & 9e-3 \\
2 um & 61 & 76\% (spots merge, FWHM 3 um) & 5.2\%, 74 nm & 7e-2 \\
\end{longtable}
}

The peak per tweezer per unit port power stays at 0.256 from 6 to 2.5 um
spacing, so power per tweezer is independent of how many are lit. At 3
um the barrier between neighbours is 78\% of the trap depth, which keeps
the atoms separate; at 2.5 um independence degrades to the 1\% level. We
take 3 um as the practical limit. It is set by the spot width (1.71 um
FWHM), hence by NA: at fixed NA, a narrower spot needs a less apodised
window, which raises the tails.

\hypertarget{3-ports-per-row-scaling-of-one-engine}{%
\section{Ports per row: scaling of one
engine}\label{scaling-3-ports-per-row-scaling-of-one-engine}}

With 1.5x oversampled arms, 2.0 um slab pitch and 1.5 um emitter pitch:

{\footnotesize
\begin{longtable}[]{@{}>{\raggedright\arraybackslash}p{0.059\dimexpr(\linewidth-20\tabcolsep)\relax}>{\raggedright\arraybackslash}p{0.059\dimexpr(\linewidth-20\tabcolsep)\relax}>{\raggedright\arraybackslash}p{0.134\dimexpr(\linewidth-20\tabcolsep)\relax}>{\raggedright\arraybackslash}p{0.064\dimexpr(\linewidth-20\tabcolsep)\relax}>{\raggedright\arraybackslash}p{0.059\dimexpr(\linewidth-20\tabcolsep)\relax}>{\raggedright\arraybackslash}p{0.081\dimexpr(\linewidth-20\tabcolsep)\relax}>{\raggedright\arraybackslash}p{0.070\dimexpr(\linewidth-20\tabcolsep)\relax}>{\raggedright\arraybackslash}p{0.157\dimexpr(\linewidth-20\tabcolsep)\relax}>{\raggedright\arraybackslash}p{0.128\dimexpr(\linewidth-20\tabcolsep)\relax}>{\raggedright\arraybackslash}p{0.192\dimexpr(\linewidth-20\tabcolsep)\relax}@{}}
\toprule\noalign{}
ports & row field & tweezers at 4 um / 3 um & slab radius & slab area &
fast actuators & static trims & data rate (200 kHz, 16 bit) & worst
aberration (rms) & full row at 4 um, 1.2 mK / 0.5 mK \\
\midrule\noalign{}
\endhead
\bottomrule\noalign{}
\endlastfoot
128 & 136 um & 30 / 40 & 1.3 mm & 0.5 mm\^{}2 & 255 & 192 & 0.8 Gbit/s &
0.03 rad & 1.3 W / 0.54 W \\
192 & 204 um & 45 / 60 & 2.0 mm & 1.1 mm\^{}2 & 383 & 288 & 1.2 Gbit/s &
0.05 rad & 1.9 W / 0.81 W \\
256 & 272 um & 60 / 80 & 2.7 mm & 2.0 mm\^{}2 & 511 & 384 & 1.6 Gbit/s &
0.07 rad & 2.6 W / 1.1 W \\
512 & 545 um & 120 / 160 & 5.3 mm & 8.2 mm\^{}2 & 1023 & 768 & 3.3
Gbit/s & 0.13 rad & 5.2 W / 2.2 W \\
1024 & 1090 um & 240 / 320 & 10.6 mm & 33 mm\^{}2 & 2047 & 1536 & 6.6
Gbit/s & 0.27 rad & 10.4 W / 4.3 W \\
\end{longtable}
}

(Row field used: 88\% of the port field. Power: 43 mW per 1.2 mK
tweezer, the D1 value, assumed to hold at larger N.)

\textbf{Aberration.} With the arm envelope produced by the synthesis,
the worst-port residual grows linearly with size, about 0.033 rad x
(ports/128), on the Rowland circle (port arc radius R/2). A quartic
correction of the port curve gains nothing, so the remainder comes from
the circular arm arc; an aspheric arm curve is the next step beyond
\textasciitilde512 ports. These geometric values are about half of the
BPM figure for D1 (0.079 rad), so they are lower bounds. Even doubled,
the Strehl ratio stays above 0.93 up to 512 ports.

\hypertarget{4-what-sets-the-limit-and-what-pushing-it-requires}{%
\section{What sets the limit, and what pushing it
requires}\label{scaling-4-what-sets-the-limit-and-what-pushing-it-requires}}

(i) \textbf{Objective field.} Each port adds one resolution cell (1.06
um at NA 0.5) to the row, so the row can only be as long as the
objective\textquotesingle s diffraction-limited field. The +/-100 um
target corresponds to about 192 ports and 60 tweezers per row. Going
further needs a larger corrected field at NA 0.5, which is the hardest
optical requirement in the system.

(ii) \textbf{Fast actuators.} Two per port: 383 for a 192-port row, each
updated at \textgreater= 200 kHz (the transport-heating limit at 1.2 mK)
and accurate to about 1\% in amplitude and 10 mrad in phase for 10 nm
position noise. Beyond a few hundred channels per chip this needs
drivers flip-chip bonded or co-integrated with the photonics, and
closed-loop calibration (on-chip power taps or camera feedback) against
PZT hysteresis and drift. This is the main engineering limit.

(iii) \textbf{Power.} About 43 mW per tweezer at 1.2 mK, 18 mW at 0.5
mK. A full 192-port row needs 1.9 W (0.8 W at 0.5 mK), which should
enter through 2 to 3 fibres feeding sub-trees, to keep each facet below
about 20 MW/cm\^{}2 (\texttt{m6\_power}). Lowering the transport depth
is the cheapest lever: it scales power linearly and relaxes the update
rate (\texttt{m3\_update\_rate}).

(iv) \textbf{Chip size.} The slab radius grows linearly with ports (2 mm
at 192, 5 mm at 512). This is within one die but beyond a single
electron-beam field, so stitching errors must be absorbed by the static
trims, or the slab written by DUV lithography.

(v) \textbf{Second dimension.} One engine addresses one row. Stacking
chips at 0.5 mm pitch, each with its own cylindrical collimator (f = 175
um, NA 0.9), and imaging the stack with 125x vertical demagnification
gives rows 4 um apart at NA 0.5. This relay is now simulated
(\texttt{docs/stacked\_relay.md}). Rows are independent provided every
chip runs at its own laser frequency; with one shared frequency,
crosstalk of about 1e-4 between rows interferes and changes neighbouring
depths by up to 2.6\%. The collimators must be aligned to about 0.1 um
in height. The cross axis and the emitter facet loss
(\texttt{docs/collimator.md}) raise the power per tweezer from 43 mW to
about 57 mW at 1.2 mK (factor 1.34) with a collimator glued to the
facets; the power column of the table above does not include this
factor. The practical limit is mechanical, set by stacking and
alignment, and plausibly 8 to 16 rows. With 192 ports per chip, 16 rows
give about 720 (4 um spacing) to 960 (3 um) independent tweezers over
about 200 x 64 um, needing about 41 to 55 W at 1.2 mK, or 17 to 23 W at
0.5 mK, spread over 16 chips (one input fibre each).

Follow-up (\texttt{docs/scaling\_beyond.md}): on a 4 um site lattice the
useful number per row is the row length / 4 um (30 for D1, 45 for 192
ports). Transport at 0.2 to 0.5 mK with 200 to 400 us moves cuts the
power 2.4 to 6 times, and a row-select tilt lets one 16-chip stack serve
all 50 rows of the +/-100 um field (2,250 sites).

\hypertarget{5-next-steps}{%
\section{Next steps}\label{scaling-5-next-steps}}

(i) Done: stacked-chip relay, \texttt{docs/stacked\_relay.md} (next:
real collimator and relay lens design). (ii) Design an aspheric arm
curve for \textgreater{} 512 ports. (iii) Specify the driver and
calibration architecture for 383 fast channels per chip. (iv) Tolerance
study of the dense port tapers (0.3 um gaps).

\hypertarget{verification-notes}{%
\section{Verification notes}\label{scaling-verification-notes}}

{\footnotesize
\begin{longtable}[]{@{}>{\raggedright\arraybackslash}p{0.403\dimexpr(\linewidth-8\tabcolsep)\relax}>{\raggedright\arraybackslash}p{0.081\dimexpr(\linewidth-8\tabcolsep)\relax}>{\raggedright\arraybackslash}p{0.114\dimexpr(\linewidth-8\tabcolsep)\relax}>{\raggedright\arraybackslash}p{0.403\dimexpr(\linewidth-8\tabcolsep)\relax}@{}}
\toprule\noalign{}
value & location & type & pointer \\
\midrule\noalign{}
\endhead
\bottomrule\noalign{}
\endlastfoot
packing table & Sec. 2 & SIM & results/m12\_packing.json (BPM matrix of
D1; review before quoting) \\
aberration vs size & Sec. 3 & SIM & results/m12\_aberration.json (exact
path lengths) \\
scaling table & Sec. 3 & DERIVED & results/m12\_table.json \\
43 mW per tweezer & Sec. 3, 4 & DERIVED & m10\_route\_E\_LS\_opt\_A.json
with the m9\_E\_power loss chain \\
crossed-AOD median 13, max 38 & Sec. 1 & USER & proposal, Table 1,
citing its ref. {[}17{]} \\
stacked-chip relay (0.5 mm pitch, 125x, 8 nm per um, factor 1.10,
frequency per chip) & Sec. 4 & SIM/DERIVED & results/m13\_*.json,
docs/stacked\_relay.md \\
8 to 16 rows & Sec. 4 & CHOICE & engineering judgement
(laser-diode-stack analogy), not a result \\
\end{longtable}
}
